\documentclass[10pt, letterpaper]{article}

% Inhaltsverzeichnis für Pakettypen (nur für Übersicht im Header, wird nicht im Dokument angezeigt)
% 1. Seitenlayout und Ränder
% 2. Sprache und Zeichensatz
% 3. Mathematik und Theorem-Umgebungen
% 4. Eigene Makros
% 5. Diagramme und Grafiken
% 6. Tabellen und Aufzählungen
% 7. Inhaltsverzeichnis
% 8. Abschnittsüberschriften
% 9. Abstrakt-Umgebung
% 10. Todos/Notizen
% 11. Rahmen/Box-Umgebungen
% 12. Python-Integration
% 13. Literaturverwaltung
% 14. Hyperlinks
% 15. Absatzeinstellungen
% 16. Umgebungen
% 17  Graphik
% 18  Extra
% 00. Titel und Autor

\usepackage{slashed}

% --- 1. Seitenlayout und Ränder ---
\usepackage[margin=3cm]{geometry}

% --- 2. Sprache und Zeichensatz ---
\usepackage[english]{babel}
\usepackage[T1]{fontenc}
\usepackage[utf8]{inputenc}

% --- 3. Mathematik und Theorem-Umgebungen ---
\usepackage{amsmath, amssymb, amsthm}
\usepackage{mathrsfs}
\DeclareMathOperator{\WF}{WF}

% --- 4. Eigene Makros ---
\usepackage{xcolor}
\newcommand{\SKP}{\langle\cdot,\cdot\rangle}
\newcommand{\id}{\operatorname{id}}
\newcommand{\sgn}{\operatorname{sgn}}
\newcommand{\Ad}{\operatorname{Ad}}
\newcommand{\ad}{\operatorname{ad}}
\newcommand{\K}{\mathbb{K}}
\newcommand{\R}{\mathbb{R}}
\newcommand{\N}{\mathbb{N}}
\newcommand{\Q}{\mathbb{Q}}
\newcommand{\Z}{\mathbb{Z}}
\newcommand{\C}{\mathbb{C}}
\newcommand{\QU}{\mathbb{H}}
\newcommand{\Cinf}{C^{\infty}}
\newcommand{\Mod}{\mathrm{Mod}}
\newcommand{\Alg}{\mathrm{Alg}}
\newcommand{\Diff}{\mathrm{Diff}}
\newcommand{\Iso}{\mathrm{Iso}}
\newcommand{\Aut}{\mathrm{Aut}}
\newcommand{\End}{\mathrm{End}}
\newcommand{\Hom}{\mathrm{Hom}}
\newcommand{\Cl}{\mathrm{Cl}}
\newcommand{\entwurf}[1]{\textcolor{red}{#1}}
\newcommand{\GL}{\mathrm{GL}}
\newcommand{\Mat}{\mathrm{Mat}}
\newcommand{\SL}{\mathrm{SL}}
\newcommand{\SO}{\mathrm{SO}}
\newcommand{\SU}{\mathrm{SU}}
\newcommand{\Sp}{\mathrm{Sp}}
\newcommand{\Spin}{\mathrm{Spin}}
\newcommand{\Pin}{\mathrm{Pin}}
\newcommand{\Lor}{\mathrm{Lor}}
\newcommand{\Ogroup}{\mathrm{O}}
\newcommand{\U}{\mathrm{U}}
\newcommand{\CP}{\mathbb{C} P}
\newcommand{\RP}{\mathbb{R} P}

% --- 5. Diagramme und Grafiken ---
\usepackage{graphicx}
\graphicspath{{../../Library/Images/}}
\usepackage[export]{adjustbox}
\usepackage{tikz}
\usetikzlibrary{decorations.pathreplacing, arrows.meta, positioning}
\usepackage{tikz-cd}

% --- 6. Tabellen und Aufzählungen ---
\usepackage{enumitem}
\setlist[itemize]{left=0.5cm}

\newenvironment{romanenum}[1][]
  {%
    \ifx&#1&
    \else
      \textbf{#1}\quad
    \fi
    \begin{enumerate}[label=\roman*)]
  }
  {%
    \end{enumerate}%
  }

% --- 7. Inhaltsverzeichnis ---
\usepackage{tocloft}
\renewcommand{\cftsecfont}{\footnotesize}
\renewcommand{\cftsubsecfont}{\footnotesize}
\renewcommand{\cftsubsubsecfont}{\footnotesize}
\renewcommand{\cftsecpagefont}{\footnotesize}
\renewcommand{\cftsubsecpagefont}{\footnotesize}
\renewcommand{\cftsubsubsecpagefont}{\footnotesize}
\usepackage{etoc}

% --- 8. Abschnittsüberschriften ---
\usepackage{titlesec}
\titleformat{\section}{\normalfont\large\bfseries}{\thesection}{1em}{}
\titleformat{\subsection}{\normalfont\normalsize\bfseries}{\thesubsection}{0.5em}{}
\titleformat{\subsubsection}{\normalfont\normalsize\bfseries}{\thesubsubsection}{0.5em}{}
\setcounter{secnumdepth}{4}

% --- 9. Abstrakt-Umgebung ---
\usepackage{changepage}
\renewenvironment{abstract}
  {
    \begin{adjustwidth}{1.5cm}{1.5cm}
    \small
    \textsc{Abstract. –}%
  }
  {
    \end{adjustwidth}
  }

% --- 10. Todos/Notizen ---
\usepackage{todonotes}
% Hellblau definieren
\definecolor{hellblau}{rgb}{0.3,0.6,1}
\newenvironment{note}{\par\begingroup\color{hellblau}\itshape}{\par\endgroup}

% --- 11. Rahmen/Box-Umgebungen ---
\usepackage{mdframed}
\usepackage{tcolorbox}
\colorlet{shadecolor}{gray!25}

\newenvironment{customTheorem}
  {\vspace{10pt}%
   \begin{mdframed}[
     backgroundcolor=gray!20,
     linewidth=0pt,
     innertopmargin=10pt,
     innerbottommargin=10pt,
     skipabove=\dimexpr\topsep+\ht\strutbox\relax,
     skipbelow=\topsep,
   ]}
  {\end{mdframed}
   \vspace{10pt}%
  }

% --- 12. Python-Integration ---
% (Deaktiviert in dieser Version, aktiviere bei Bedarf)
% \usepackage{pythontex}
% \usepackage[makestderr]{pythontex}

% --- 13. Literaturverwaltung ---
\usepackage{csquotes}
\usepackage[backend=biber, style=alphabetic, citestyle=alphabetic]{biblatex}
\addbibresource{bibliography.bib}

% --- 14. Hyperlinks ---
\usepackage{hyperref}
\hypersetup{
  colorlinks   = true,
  urlcolor     = blue,
  linkcolor    = blue,
  citecolor    = blue,
  frenchlinks  = true
}

% --- 15. Absatzeinstellungen ---
\usepackage[parfill]{parskip}
\sloppy

% --- 16. Umgebungen ---
\usepackage{thmtools}

\newcommand{\CustomHeading}[3]{%
  \par\medskip\noindent%
  \textbf{#1 #2} \textnormal{(#3)}.\enskip%
}

\newenvironment{DEF}[2]{\CustomHeading{Definition}{#1}{#2}}{}
\newenvironment{PROP}[2]{\CustomHeading{Proposition}{#1}{#2}}{}
\newenvironment{THEO}[2]{\CustomHeading{Theorem}{#1}{#2}}{}
\newenvironment{LEM}[2]{\CustomHeading{Lemma}{#1}{#2}}{}
\newenvironment{KORO}[2]{\CustomHeading{Corollar}{#1}{#2}}{}
\newenvironment{REM}[2]{\CustomHeading{Remark}{#1}{#2}}{}
\newenvironment{EXA}[2]{\CustomHeading{Example}{#1}{#2}}{}
\newenvironment{STUD}[2]{\CustomHeading{Study}{#1}{#2}}{}
\newenvironment{CONC}[2]{\CustomHeading{Concept}{#1}{#2}}{}
\newenvironment{OTH}[2]{\CustomHeading{Other}{#1}{#2}}{}
\newenvironment{EXE}[2]{\CustomHeading{Exercise}{#1}{#2}}{}
\newenvironment{MOT}[2]{\CustomHeading{Motivation}{#1}{#2}}{}
\newenvironment{PROOF}[2]{\CustomHeading{Proof}{#1}{#2}}{}

% --- Unit Umgebung für Source-Inhalte ---
\usepackage{mdframed}
\newmdenv[
  linewidth=1pt,
  topline=false,
  bottomline=false,
  rightline=false,
  leftmargin=0cm,
  rightmargin=0cm,
  skipabove=10pt,
  skipbelow=10pt,
  innertopmargin=0.5\baselineskip,
  innerbottommargin=0.5\baselineskip,
  backgroundcolor=gray!10,
  linecolor=gray
]{unitbox}

\newenvironment{unit}[1]
  {\begin{unitbox}\textbf{Unit #1}\par\smallskip}
  {\end{unitbox}}

% --- 17. ... ---

% --- 18. Extras ---
\usepackage{stmaryrd}
\usepackage{bbold}  % falls du \mathbb{1} nutzen willst

% --- 00. Titel und Autor ---
\title{Mein Titel}
\author{Tim Jaschik}
\date{\today}

\begin{document}

\maketitle
\rule{\textwidth}{0.5pt}
\begin{abstract}
Kurze Beschreibung …
\end{abstract}
\rule{\textwidth}{0.5pt}
\vspace{0.5cm}

\tableofcontents

\pagebreak

\subsubsection{Die Spin-Gruppe als gerade Untergruppe der Clifford-Einheiten}

Sei $(\R^n,\langle\cdot,\cdot\rangle)$ ein euklidischer Vektorraum und 
$\Cl_n$ die zugehörige Clifford-Algebra.

\begin{DEF}{}{Pin- und Spin-Gruppe}
\begin{enumerate}
\item Die \emph{Pin-Gruppe} ist definiert als die von den Einheitsvektoren erzeugte Untergruppe
\[
\Pin(n):=\langle S^{n-1}\rangle \subset \Cl_n^\times.
\]

\item Die \emph{Spin-Gruppe} ist der gerade Teil davon:
\[
\Spin(n):=\Pin(n)\cap \Cl_n^+.
\]
\end{enumerate}
\end{DEF}

\begin{DEF}{}{Adjungierte Wirkung}
Sei $\alpha:\Cl_n\to\Cl_n$ die Gradinvolution, d.h.
\[
\alpha(v)=-v \quad \text{für } v\in \R^n,
\qquad
\alpha(x)=(-1)^k x \text{ für } x\in \Cl_n^{(k)}.
\]
Dann definiert man für $u\in \Pin(n)$ eine lineare Abbildung
\[
\rho(u):\R^n\to \R^n,\qquad
\rho(u)(x):=\alpha(u)\,x\,u^{-1}.
\]
\end{DEF}

\begin{PROP}{}{Geometrische Bedeutung von $\rho$}
\begin{enumerate}
\item Für $u\in S^{n-1}$ ist $\rho(u)$ die Spiegelung an der Hyperebene orthogonal zu $u$.
\item Für allgemeines $u\in \Pin(n)$ ist $\rho(u)$ eine orthogonale Transformation.
\item Es gilt:
\[
\rho:\Pin(n)\to \mathrm{O}(n)
\quad \text{ist ein surjektiver Gruppenhomomorphismus mit } \ker\rho=\{\pm 1\}.
\]
\end{enumerate}
\end{PROP}

\begin{PROP}{}{Einschränkung auf die Spin-Gruppe}
Für $u\in \Spin(n)\subset \Cl_n^+$ gilt
\[
\alpha(u)=u.
\]
Damit vereinfacht sich die Wirkung zu
\[
\rho(u)(x)=u\,x\,u^{-1}.
\]
\end{PROP}

\begin{KORO}{}{Doppelüberlagerung von $\SO(n)$}
Die Einschränkung von $\rho$ liefert einen surjektiven Gruppenhomomorphismus
\[
\rho:\Spin(n)\to \SO(n),
\qquad
\rho(u)(x)=u x u^{-1},
\]
mit
\[
\ker\rho=\{\pm 1\}.
\]
Insbesondere ist $\Spin(n)$ eine zweifache Überlagerung von $\SO(n)$.
\end{KORO}

\begin{REM}{}{Konzeptuelle Interpretation}
\begin{enumerate}
\item Die Verwendung von $\alpha$ ist notwendig, um auf ganz $\Pin(n)$ eine orthogonale Wirkung zu erhalten (insbesondere für Spiegelungen).
\item Auf $\Spin(n)$ wird $\alpha$ trivial, da alle Elemente gerade sind.
\item Die Wirkung
\[
x \mapsto u x u^{-1}
\]
ist daher die intrinsische geometrische Realisierung von Rotationen in der Clifford-Algebra.
\end{enumerate}
\end{REM}



\pagebreak

\subsubsection{Kompatibilität von Spin-Wirkung und Clifford-Multiplikation}

Sei $(M,g)$ eine orientierte Riemannsche Mannigfaltigkeit mit Spin-Struktur
\[
P^{\Spin}\xrightarrow{\;\pi\;} M.
\]

Sei ferner
\[
\rho:\Spin(n)\to \Aut(\Sigma_n)
\]
die Spinor-Darstellung und
\[
\lambda:\Spin(n)\to \SO(n)
\]
die kanonische Überlagerung.

\medskip

\textbf{Ziel:} Konstruktion einer globalen Clifford-Multiplikation
\[
c:TM\longrightarrow \End(\Sigma M).
\]

\paragraph{Lokale Definition}

Für $q\in P^{\Spin}_x$ sei
\[
u_q:\R^n \xrightarrow{\sim} T_xM
\]
die durch den Rahmen gegebene Isometrie.

Dann definieren wir lokal für $v\in T_xM$:
\[
c_x(v):=c_0\bigl(u_q^{-1}(v)\bigr),
\]
wobei
\[
c_0:\R^n\to \End(\Sigma_n)
\]
die Standard-Clifford-Wirkung ist.

\paragraph{Problem: Abhängigkeit von der Wahl von $q$}

Für $g\in \Spin(n)$ gilt:
\[
q' = q\cdot g
\quad \Longrightarrow \quad
u_{q'} = u_q \circ \lambda(g).
\]

Damit folgt:
\[
u_{q'}^{-1}(v)
=
\lambda(g)^{-1}\,u_q^{-1}(v).
\]

Also ergibt sich für die alternative Wahl:
\[
c_x'(v)
=
c_0\bigl(\lambda(g)^{-1} u_q^{-1}(v)\bigr).
\]

\paragraph{Transformationsverhalten der Spinoren}

Ein Spinorfeld transformiert unter Rahmenwechsel durch
\[
\psi \longmapsto \rho(g)\psi.
\]

\paragraph{Konsistenzforderung}

Damit die Clifford-Multiplikation wohldefiniert ist, muss gelten:
\[
c_x'(v)\,\rho(g)\psi
=
\rho(g)\,c_x(v)\psi.
\]

Einsetzen liefert:
\[
c_0(\lambda(g)^{-1} w)\,\rho(g)
=
\rho(g)\,c_0(w),
\quad w:=u_q^{-1}(v).
\]

Äquivalent:
\[
\rho(g)\,c_0(w)\,\rho(g)^{-1}
=
c_0(\lambda(g) w).
\]

\paragraph{Zentrale Kompatibilitätsbedingung}

\begin{equation}
\boxed{
c_0(\lambda(g) v)
=
\rho(g)\,c_0(v)\,\rho(g)^{-1}
}
\end{equation}

\paragraph{Interpretation}

\begin{enumerate}
\item Die Spin-Darstellung $\rho$ ist so gewählt, dass sie die
\emph{Clifford-Multiplikation äquivariant} macht.
\item Die Wirkung von $\Spin(n)$ auf $\R^n$ über $\lambda$ entspricht genau der
\emph{Konjugationswirkung} in der Clifford-Algebra:
\[
\lambda(g)v = g v g^{-1}.
\]
\item Die Relation
\[
\rho(g)(v\cdot \psi)
=
(\lambda(g)v)\cdot (\rho(g)\psi)
\]
ist genau die Bündelversion dieser Kompatibilität.
\end{enumerate}

\paragraph{Schluss}

Diese Bedingung ist notwendig und hinreichend dafür, dass die lokal definierte Abbildung
\[
c_x(v):=c_0(u_q^{-1}(v))
\]
\emph{unabhängig von der Wahl des Spin-Rahmens} ist und somit eine globale Bündelabbildung
\[
c:TM\to \End(\Sigma M)
\]
definiert.


\pagebreak

\begin{THEO}{}{Äquivarianz der Clifford-Multiplikation unter der Spin-Wirkung}
Sei $\Cl_n$ die Clifford-Algebra über $(\R^n,\langle\cdot,\cdot\rangle)$,
$\Sigma_n$ ein $\Cl_n$-Modul und
\[
\rho:\Spin(n)\to \Aut(\Sigma_n)
\]
die durch Einschränkung der Clifford-Wirkung definierte Darstellung,
d.h.
\[
\rho(g)\psi := g\cdot \psi.
\]

Dann gilt für alle $g\in \Spin(n)$ und $v\in \R^n$:
\[
\rho(g)\,c_0(v)\,\rho(g)^{-1}
=
c_0(\lambda(g)v),
\]
wobei
\[
\lambda(g)v := g v g^{-1}.
\]
\end{THEO}

\begin{proof}
Sei $\psi\in \Sigma_n$.

Dann gilt zunächst nach Definition der Clifford-Multiplikation:
\[
c_0(v)\psi = v\cdot \psi.
\]

Weiter ist
\[
\rho(g)\psi = g\cdot \psi,
\qquad
\rho(g)^{-1}\psi = g^{-1}\cdot \psi.
\]

Damit ergibt sich:
\[
\begin{aligned}
\rho(g)\,c_0(v)\,\rho(g)^{-1}(\psi)
&=
\rho(g)\bigl(v\cdot (g^{-1}\cdot \psi)\bigr) \\
&=
g\cdot \bigl(v\cdot (g^{-1}\cdot \psi)\bigr).
\end{aligned}
\]

Da $\Sigma_n$ ein Modul über der assoziativen Algebra $\Cl_n$ ist,
gilt die Assoziativität:
\[
g\cdot (v\cdot \phi) = (g v)\cdot \phi.
\]

Somit folgt:
\[
\begin{aligned}
g\cdot \bigl(v\cdot (g^{-1}\cdot \psi)\bigr)
&=
(g v)\cdot (g^{-1}\cdot \psi) \\
&=
(g v g^{-1})\cdot \psi.
\end{aligned}
\]

Nach Definition von $\lambda$ ist
\[
\lambda(g)v = g v g^{-1}.
\]

Also erhalten wir:
\[
\rho(g)\,c_0(v)\,\rho(g)^{-1}(\psi)
=
(\lambda(g)v)\cdot \psi
=
c_0(\lambda(g)v)\psi.
\]

Da dies für alle $\psi\in \Sigma_n$ gilt, folgt die Behauptung.
\end{proof}

\pagebreak
\begin{THEO}{}{Existenz der globalen Clifford-Multiplikation auf dem Spinorbündel}
Sei $(M,g)$ eine orientierte Riemannsche Spin-Mannigfaltigkeit mit Spin-Struktur
\[
\pi: P^{\Spin}(M) \longrightarrow M
\]
und zugehörigem Spinorbündel
\[
\Sigma M := P^{\Spin}(M)\times_{\rho} \Sigma_n,
\]
wobei $\rho:\Spin(n)\to \Aut(\Sigma_n)$ die Spin-Darstellung ist.
Dann existiert eine eindeutig bestimmte glatte Bündelabbildung
\[
c: TM \longrightarrow \End(\Sigma M),
\]
so dass für jedes $x\in M$ eine Darstellung der Clifford-Algebra
\[
c_x: \Cl(T_xM,g_x)\longrightarrow \End(\Sigma_x M)
\]
induziert wird. Insbesondere gilt punktweise die Clifford-Relation
\[
c_x(v)c_x(w)+c_x(w)c_x(v) = -2g_x(v,w)\id_{\Sigma_x M}.
\]
\end{THEO}

\begin{proof}
Wir verwenden die Darstellung des Tangentialbündels und des Spinorbündels als assoziierte Bündel:
\[
TM \cong P^{\Spin}\times_{\lambda} \R^n,
\qquad
\Sigma M \cong P^{\Spin}\times_{\rho} \Sigma_n,
\]
wobei $\lambda:\Spin(n)\to \SO(n)$ die Standardwirkung ist.

\medskip

\noindent
\textbf{1. Definition.}
Sei $x\in M$, $q\in P^{\Spin}_x$, $v\in T_xM$ und $\psi\in \Sigma_n$.
Über die durch $q$ induzierte Identifikation
\[
u_q:\R^n \xrightarrow{\sim} T_xM
\]
definieren wir
\[
c_x(v)[q,\psi] := [q,\,u_q^{-1}(v)\cdot \psi],
\]
wobei die rechte Seite die Clifford-Multiplikation im Modellraum $\Sigma_n$ bezeichnet.

\medskip

\noindent
\textbf{2. Wohldefiniertheit.}
Wir müssen zeigen, dass die Definition unabhängig von der Wahl des Repräsentanten ist.

Sei $g\in \Spin(n)$. Dann gilt
\[
(q,\psi) \sim (qg,\rho(g^{-1})\psi).
\]
Zugleich transformiert die Identifikation nach
\[
u_{qg} = u_q \circ \lambda(g).
\]
Daher folgt
\[
u_{qg}^{-1}(v) = \lambda(g)^{-1} u_q^{-1}(v).
\]

Somit ergibt sich
\[
\begin{aligned}
c_x(v)[qg,\rho(g^{-1})\psi]
&= [qg,\,\lambda(g)^{-1}u_q^{-1}(v)\cdot \rho(g^{-1})\psi].
\end{aligned}
\]

Nun verwenden wir die $\Spin(n)$-Äquivarianz der Clifford-Multiplikation
\[
\rho(g^{-1})(v\cdot \psi)
=
(\lambda(g^{-1})v)\cdot (\rho(g^{-1})\psi),
\]
um zu erhalten
\[
\lambda(g)^{-1}u_q^{-1}(v)\cdot \rho(g^{-1})\psi
=
\rho(g^{-1})\bigl(u_q^{-1}(v)\cdot \psi\bigr).
\]

Damit folgt
\[
\begin{aligned}
c_x(v)[qg,\rho(g^{-1})\psi]
&=
[qg,\,\rho(g^{-1})(u_q^{-1}(v)\cdot \psi)] \\
&=
[q,\,u_q^{-1}(v)\cdot \psi].
\end{aligned}
\]

Also ist $c_x(v)$ wohldefiniert.

\medskip

\noindent
\textbf{3. Glattheit und Linearität.}
Die Konstruktion ist lokal durch glatte Funktionen in Trivialisierungen gegeben und daher glatt.
Die Linearität in $v$ folgt unmittelbar aus der Linearität der Clifford-Multiplikation im Modellraum.

\medskip

\noindent
\textbf{4. Clifford-Relation.}
Für $v,w\in T_xM$ gilt
\[
\begin{aligned}
c_x(v)c_x(w)[q,\psi]
&=
[q,\,u_q^{-1}(v)\cdot (u_q^{-1}(w)\cdot \psi)].
\end{aligned}
\]
Da im Modellraum $\Sigma_n$ die Clifford-Relation gilt,
\[
u_q^{-1}(v)\cdot u_q^{-1}(w)
+
u_q^{-1}(w)\cdot u_q^{-1}(v)
=
-2g_x(v,w),
\]
folgt unmittelbar
\[
c_x(v)c_x(w)+c_x(w)c_x(v)
=
-2g_x(v,w)\id.
\]

\medskip

\noindent
\textbf{5. Eindeutigkeit.}
Die Abbildung ist durch die Forderung bestimmt, lokal mit der Modell-Clifford-Multiplikation übereinzustimmen. Daher ist sie eindeutig.

\medskip

Damit ist die globale Clifford-Multiplikation konstruiert.
\end{proof}


\pagebreak
\printbibliography
\end{document}